% =============================================================================
% Chapter 3: ML Workflow — ML Systems Lecture Slides
% =============================================================================
% This deck covers the ML lifecycle, workflow engineering, and systems thinking.
%
% IMAGES NEEDED (generate SVGs, convert to PDF):
%   - ch03-day150-timeline.pdf     - ch03-dual-pipeline.pdf
%   - ch03-time-allocation.pdf     - ch03-constraint-propagation.pdf
%   - ch03-iteration-velocity.pdf  - ch03-iron-law-workflow.pdf
%   - ch03-lifecycle-circle.pdf    - ch03-systems-thinking-triad.pdf
% =============================================================================
\documentclass[aspectratio=169, 12pt]{beamer}
\usepackage{../../assets/beamerthememlsys}

\mlsyssetup{
  volume   = {Volume I},
  chapter  = {Chapter 3},
  logo = {../../assets/img/logo-mlsysbook.png},
  instlogo = {../../assets/img/logo-harvard.png},
  chaptertitle = {ML Workflow},
}

% --- Fonts ---
\usepackage{fontspec}
\setsansfont{Helvetica Neue}[
  BoldFont={Helvetica Neue Bold},
  ItalicFont={Helvetica Neue Italic},
  BoldItalicFont={Helvetica Neue Bold Italic},
]
\setmonofont{JetBrains Mono}[Scale=0.85]

% --- Packages ---
\usepackage{booktabs}
\usepackage{amsmath}

% --- Image paths ---
\graphicspath{
  {images/}
}

% --- Chapter-specific macros ---
\newcommand{\DAM}{D\raisebox{0.08em}{\tiny$\bullet$}A\raisebox{0.08em}{\tiny$\bullet$}M}

% --- Helper: safe image include ---
\newcommand{\safeimg}[2][width=\textwidth,keepaspectratio]{%
  \IfFileExists{images/#2}{\includegraphics[#1]{#2}}{%
    \IfFileExists{#2}{\includegraphics[#1]{#2}}{%
      \fbox{\parbox[c][2.5cm][c]{0.85\linewidth}{\centering\footnotesize\textcolor{midgray}{[Missing image]}}}%
    }%
  }%
}

% --- Section count for navigation (must match actual \section{} count) ---
\setcounter{mlsystotalsections}{9}

\title{ML Workflow}
\author{Vijay Janapa Reddi}
\institute{Harvard University}
\date{}

\begin{document}

% =============================================================================
% TITLE SLIDE
% =============================================================================
\mlsystitle{ML Workflow}{See the Whole Map Before Walking Any Path}{cover_ml_workflow.png}

% =============================================================================
% LEARNING OBJECTIVES
% =============================================================================
\begin{frame}{Learning Objectives}
\note{[2 min] Walk through objectives. Emphasize that this chapter provides the
mental map for the entire course. Ask: ``How many of you have built a model but
never thought about how it would be deployed?''}

\small
\begin{enumerate}
  \item Describe the six \textbf{ML lifecycle stages} and their feedback-driven relationships
  \item Compare ML workflows to \textbf{traditional software} development
  \item Apply the \textbf{Constraint Propagation Principle} ($2^{N-1}$ cost escalation)
  \item Map lifecycle stages to \textbf{Iron Law} terms ($D_{\text{vol}}/BW$, $O/R_p\eta$, $L_{\text{lat}}$)
  \item Explain why \textbf{iteration velocity} compounds into quality advantages
  \item Trace how \textbf{multi-scale feedback} loops enable continuous improvement
  \item Apply \textbf{systems thinking} to anticipate cross-stage failures
\end{enumerate}

\end{frame}

% =============================================================================
\section{The Day-150 Disaster}
% =============================================================================

\begin{frame}{The Day-150 Disaster}
\note{[3 min] Open with this visceral story. Build suspense: the model was
excellent. The team was excellent. But 5 months of work was discarded.
Ask: ``What went wrong?'' Let students guess before revealing.}

% --- Layout: FULL-WIDTH timeline ---
\centering
\safeimg[width=0.92\textwidth,height=4.5cm,keepaspectratio]{ch03-day150-timeline.pdf}

\vspace{0.15cm}
\mlsysalert{Key:}{The model's accuracy was excellent. The \emph{workflow} failed.}

\end{frame}

\begin{frame}{What Went Wrong?}
\note{[2 min] Break down the failure. The tablet's 512 MB memory limit should
have shaped every decision from Day 1. Instead, the team optimized each
component in isolation. This motivates the entire chapter.}

\footnotesize
\begin{columns}[T]
  \begin{column}{0.52\textwidth}
    \textbf{The timeline:}
    \begin{itemize}\setlength\itemsep{0pt}
      \item Day 1: ``Build a diagnostic model''
      \item Day 90: 95\% accuracy on test set
      \item Day 120: 96\% after architecture tuning
      \item Day 150: Handed to deployment
      \item Day 151: Model requires \textbf{4 GB} memory
      \item Day 152: Tablet has \textbf{512 MB}
      \item Day 153: Five months discarded
    \end{itemize}
  \end{column}
  \begin{column}{0.44\textwidth}
    \begin{mlsyscard}{errorstroke}
    \textbf{Root cause:}\\[0.05cm]
    {\scriptsize A deployment constraint ($8\times$ memory violation) that should have shaped Day 1 was discovered on Day 151.}
    \end{mlsyscard}

    \vspace{0.1cm}
    \begin{mlsyscard}{crimson}
    {\scriptsize This chapter introduces the \textbf{ML Workflow}: engineering that prevents such failures.}
    \end{mlsyscard}
  \end{column}
\end{columns}

\end{frame}

\begin{frame}{The ML Lifecycle: Definition}
\note{[2 min] Formal definition. Emphasize two key distinctions: (1) lifecycle
is WHAT gets traversed, workflow is HOW. (2) Unlike traditional software,
ML systems degrade through data drift even when code is untouched.
Ask: ``What degrades in ML but not in traditional software?''}

\footnotesize
\begin{mlsyscard}{crimson}
\textbf{Machine Learning Lifecycle:} The continuous engineering discipline of managing
\textbf{System Entropy} across the Data, Algorithm, and Machine axes.
\end{mlsyscard}

\vspace{0.1cm}
\renewcommand{\arraystretch}{1.1}
\begin{tabular}{@{}lll@{}}
  \toprule
  & \textbf{Traditional SW} & \textbf{ML System} \\
  \midrule
  Degrades via    & Code modification      & Data drift (untouched code) \\
  Release model   & Ship and maintain      & Continuous retraining loop \\
  End of lifecycle & Deployment             & \textbf{Beginning} of feedback loop \\
  \bottomrule
\end{tabular}

\vspace{0.1cm}
\mlsysconcept{Key distinction:}{\textit{Lifecycle} = stages traversed. \textit{Workflow} = how traversal is managed.}

\end{frame}

% =============================================================================
\section{ML vs.\ Software}
% =============================================================================

\begin{frame}{Where Time Actually Goes}
\note{[3 min] This slide shocks students. Most expect model development to
dominate. Show the stacked bar: 79\% of effort is data-related. Only 4\%
goes to refining algorithms. Ask: ``Where would you invest your next hour?''}

% --- Layout: FULL-WIDTH image ---
\centering
\safeimg[width=0.92\textwidth,height=4.8cm,keepaspectratio]{ch03-time-allocation.pdf}

\vspace{0.1cm}
\mlsysinsight{Implication:}{Investment in data engineering yields the highest returns per hour spent.}

\end{frame}

\begin{frame}{ML vs.\ Traditional Software}
\note{[3 min] Walk through six dimensions of difference. The most critical is
the last row: feedback loops. Traditional software rarely has later stages
influence earlier ones. ML systems require it. If short: focus on feedback
loops and failure modes.}

\scriptsize
\renewcommand{\arraystretch}{1.1}
\begin{tabular}{@{}lp{3.5cm}p{5cm}@{}}
  \toprule
  \textbf{Dimension} & \textbf{Traditional Software} & \textbf{ML Systems} \\
  \midrule
  \textbf{Problem Def}  & Precise specs upfront         & Objectives evolve through exploration \\
  \textbf{Development}   & Linear feature implementation & Iterative experimentation with data \\
  \textbf{Testing}       & Deterministic pass/fail       & Statistical validation with uncertainty \\
  \textbf{Deployment}    & Static until updated          & Drifts as data distributions shift \\
  \textbf{Maintenance}   & Fix bugs, add features        & Continuous monitoring + retraining \\
  \textbf{Feedback}      & Minimal; later $\not\to$ earlier & \textcolor{errorstroke}{\textbf{Frequent; deploy $\to$ data}} \\
  \bottomrule
\end{tabular}

\vspace{0.1cm}
\small
\mlsysalert{Core difference:}{ML systems require continuous feedback loops that traditional software never needs.}

\end{frame}

\begin{frame}{The Iteration Tax}
\note{[3 min] Worked example. Two models compete over 6 months. The large
model starts better but iterates slowly. The small model starts 5\% behind
but iterates 168$\times$ faster. Ask: ``Which wins?'' Reveal after pause.}

\footnotesize
\begin{columns}[T]
  \begin{column}{0.48\textwidth}
    \safeimg[width=\textwidth,height=5.0cm,keepaspectratio]{ch03-iteration-velocity.pdf}
  \end{column}
  \begin{column}{0.48\textwidth}
    \textbf{6-month competition:}

    \vspace{0.05cm}
    \textcolor{computestroke}{\textbf{Large Model}}\\
    {\scriptsize Start: 95\%, cycle: 1 week, +0.15\%/iter. 26 iters $\to$ \textbf{99\%}}

    \vspace{0.05cm}
    \textcolor{datastroke}{\textbf{Small Model}}\\
    {\scriptsize Start: 90\%, cycle: 1 hour, +0.1\%/iter. 100 iters $\to$ \textbf{99\%}}

    \vspace{0.1cm}
    \begin{mlsyscard}{datastroke}
    {\scriptsize \textbf{Iteration velocity is a feature.}\\
    10 experiments/day almost always outperforms 1/week.}
    \end{mlsyscard}
  \end{column}
\end{columns}

\end{frame}

% --- ACTIVE LEARNING 1: Predict ---
\begin{frame}{Predict: What Happens Without Problem Definition?}
\note{[2 min] Prediction exercise. Give students 60 seconds. Do NOT reveal
the answer. This primes the Constraint Propagation Principle.
Ask 2--3 students to share answers.}

\centering
\vspace{0.6cm}
{\Large\bfseries Think--Write--Share}

\vspace{0.4cm}
{\large A team builds a highly accurate model (98\%)\\
but skips the Problem Definition stage entirely.\\[0.2cm]
\alert{What goes wrong at deployment?}}

\vspace{0.4cm}
{\normalsize Write down 2 specific failures. \textcolor{midgray}{(60 seconds)}}

\vspace{0.3cm}
{\small\textcolor{lightgray}{Turn to a neighbor and compare your answers.}}

\end{frame}

% =============================================================================
\section{Lifecycle Stages}
% =============================================================================

\begin{frame}{The Six Lifecycle Stages}
\note{[3 min] Overview diagram. Emphasize the feedback loop from Monitoring
back to Problem Definition. This is NOT a waterfall. The misconception that
``lifecycle ends at deployment'' is the most dangerous one.
If short: just cover the loop, not all 6 stages.}

% --- Layout: FULL-WIDTH diagram ---
\centering
\safeimg[width=0.88\textwidth,height=4.5cm,keepaspectratio]{ch03-lifecycle-circle.pdf}

\vspace{0.15cm}
\mlsysinsight{Key:}{Deployment is the \emph{beginning} of the feedback loop, not the end.}

\end{frame}

\begin{frame}{The Dual-Pipeline Architecture}
\note{[3 min] Two parallel pipelines running simultaneously. Data pipeline
(green, top) and Model pipeline (blue, bottom). The backward feedback arrows
are what make ML systems fundamentally different from traditional software.
Ask: ``Which pipeline gets more research attention? Which consumes more time?''}

% --- Layout: FULL-WIDTH diagram ---
\centering
\safeimg[width=0.92\textwidth,height=4.5cm,keepaspectratio]{ch03-dual-pipeline.pdf}

\vspace{0.1cm}
\footnotesize
\textbf{Two pipelines, one system.} Data and model development run in parallel,\\
unified by continuous feedback loops that have no equivalent in traditional software.

\end{frame}

\begin{frame}{Case Study: DR Screening}
\note{[2 min] Introduce the running case study. Diabetic retinopathy screening
seems simple (classify retinal images) but reveals every workflow challenge.
Key facts: 100M+ affected worldwide, 1 ophthalmologist per 100K+ people in
rural areas. AI-assisted screening is medically essential.}

\footnotesize
\begin{columns}[T]
  \begin{column}{0.52\textwidth}
    \textbf{Diabetic Retinopathy Screening:}
    \begin{itemize}\setlength\itemsep{0pt}
      \item Affects 100M+ people worldwide
      \item Leading cause of preventable blindness
      \item 1 ophthalmologist per 100K+ in rural areas
      \item AI screening: medically essential
    \end{itemize}

    \vspace{0.05cm}
    \textbf{Why this case study?}
    \begin{itemize}\setlength\itemsep{0pt}
      \item Simple surface, deep complexity
      \item Spans the deployment spectrum
      \item Well-documented journey to production
    \end{itemize}
  \end{column}
  \begin{column}{0.44\textwidth}
    \begin{mlsyscard}{crimson}
    \textbf{Five competing constraints:}\\[0.05cm]
    {\scriptsize
    1. Diagnostic accuracy\\
    2. Computational efficiency\\
    3. Workflow integration\\
    4. Regulatory compliance\\
    5. Cost-effectiveness}
    \end{mlsyscard}

    \vspace{0.05cm}
    {\scriptsize\textcolor{midgray}{Each constraint tightens the feasible design space for the others.}}
  \end{column}
\end{columns}

\end{frame}

% =============================================================================
\section{Iron Law of Workflow}
% =============================================================================

\mlsysfocus{The Iron Law of Workflow}{%
Each lifecycle stage maps to a specific term in:\\[0.3cm]
$T_{\text{total}} = \underbrace{\dfrac{D_{\text{vol}}}{BW}}_{\text{Data Collection}} +
\underbrace{\dfrac{O}{R_{\text{peak}} \cdot \eta}}_{\text{Model Dev}} +
\underbrace{L_{\text{lat}}}_{\text{Deployment}}$%
}

\begin{frame}{Iron Law of Workflow: Stages Map to Terms}
\note{[3 min] Bridge the workflow to the Iron Law from Ch1. Each lifecycle
stage optimizes a specific term. Data Collection determines D\_vol.
Model Development defines O. Deployment minimizes L\_lat.
Ask: ``If you only optimize the model (compute term), what happens?''}

% --- Layout: FULL-WIDTH diagram ---
\centering
\safeimg[width=0.92\textwidth,height=4.8cm,keepaspectratio]{ch03-iron-law-workflow.pdf}

\vspace{0.1cm}
\small
\mlsysconcept{Key insight:}{Workflow management is mathematically equivalent to minimizing the Iron Law.}

\end{frame}

\begin{frame}{Lighthouse Models: Different Workflows}
\note{[3 min] Same lifecycle stages, different dominant Iron Law terms.
ResNet-50 is compute-bound. DLRM is I/O-bound. KWS is energy-bound.
Each requires fundamentally different workflow priorities.
If short: just contrast ResNet-50 vs KWS.}

\scriptsize
\renewcommand{\arraystretch}{1.05}
\begin{tabular}{@{}lp{3.2cm}p{3.2cm}p{3.2cm}@{}}
  \toprule
  \textbf{Stage} & \textbf{ResNet-50} \newline\textcolor{computestroke}{Compute Beast} & \textbf{DLRM} \newline\textcolor{datastroke}{Sparse Scatter} & \textbf{KWS} \newline\textcolor{errorstroke}{Tiny Constraint} \\
  \midrule
  \textbf{Data}    & Throughput: >80\% GPU utilization  & Feature store: <2 ms lookups & Fit in 256 KB RAM \\
  \textbf{Train}   & Maximize $\eta$; mixed precision   & Optimize sparse embeddings   & NAS for smallest model \\
  \textbf{Deploy}  & Batch >128 for throughput          & Strict <10 ms p99 latency    & <1 mW energy budget \\
  \bottomrule
\end{tabular}

\vspace{0.1cm}
\small
\mlsysconcept{Key:}{Same workflow stages, different Iron Law terms dominate by archetype.}

\end{frame}

% =============================================================================
\section{Data \& Models}
% =============================================================================

\begin{frame}{Stage 2: Data Collection Challenges}
\note{[3 min] Data collection is the primary engineering activity. For DR
screening: 128K retinal images, multiple expert annotations, diverse clinic
equipment. The bandwidth constraint forces edge deployment.
Common error: students think ``just get more data.''}

\footnotesize
\begin{columns}[T]
  \begin{column}{0.52\textwidth}
    \textbf{DR Screening Data Challenges:}
    \begin{itemize}\setlength\itemsep{0pt}
      \item $10^5$ retinal photographs needed
      \item Multiple expert ophthalmologist reviews
      \item Diverse cameras, lighting, operators
      \item Rural clinics: limited connectivity
    \end{itemize}

    \vspace{0.05cm}
    \textbf{Bandwidth vs.\ Compute:}\\
    {\scriptsize 150 patients $\times$ 10 photos $\times$ 5 MB = \textbf{7 GB/day}\\
    At 2 Mbps = \textbf{25+ hours} to upload (clinic: 8 hours)}
  \end{column}
  \begin{column}{0.44\textwidth}
    \begin{mlsyscard}{datastroke}
    {\scriptsize \textbf{Engineering conclusion:} Cloud-only is bandwidth-infeasible. Edge reduces uploads by \textbf{5,000$\times$}.}
    \end{mlsyscard}

    \vspace{0.1cm}
    \begin{mlsyscard}{routingstroke}
    {\scriptsize \textbf{Lab-to-Field Gap:} Models trained on research images fail on blurry, poorly-lit clinic images.}
    \end{mlsyscard}
  \end{column}
\end{columns}

\end{frame}

\begin{frame}{Stage 3: Model Development}
\note{[3 min] The DR team has 128K images and a target: >90\% sensitivity on
edge hardware with <50 ms latency. Key techniques: transfer learning (reuse
ImageNet representations), model compression. The competition-production gap:
Netflix Prize ensemble won \$1M but was never deployed.}

\footnotesize
\begin{columns}[T]
  \begin{column}{0.52\textwidth}
    \textbf{DR Model Development:}
    \begin{itemize}\setlength\itemsep{0pt}
      \item \textbf{Transfer learning}: Reuse ImageNet, fine-tune with domain data
      \item Result: AUC 0.99, sensitivity 97.5\%
      \item But: initial model too large for edge
    \end{itemize}

    \vspace{0.05cm}
    \textbf{Accuracy vs.\ Efficiency:}
    \begin{itemize}\setlength\itemsep{0pt}
      \item Quantization: FP32 $\to$ INT8
      \item Pruning: remove redundant connections
      \item Distillation: large $\to$ small model
    \end{itemize}
  \end{column}
  \begin{column}{0.44\textwidth}
    \begin{mlsyscard}{errorstroke}
    {\scriptsize \textbf{Competition-Production Gap:} Netflix Prize: 800+ model ensemble. \textbf{Never deployed} --- serving complexity exceeded value.}
    \end{mlsyscard}

    \vspace{0.05cm}
    \begin{mlsyscard}{computestroke}
    {\scriptsize \textbf{Reproducible artifact:} Weights + code + environment + config. ``Works on my machine'' = broken system.}
    \end{mlsyscard}
  \end{column}
\end{columns}

\end{frame}

\begin{frame}{Stage Interface Contracts}
\note{[2 min] Each lifecycle stage has explicit inputs, outputs, and quality
invariants. Think of these as API contracts between teams. When a stage
output fails its contract, costs compound downstream. This is the formal
mechanism that prevents Day-150 disasters.}

\scriptsize
\renewcommand{\arraystretch}{1.05}
\begin{tabular}{@{}lp{3cm}p{3.5cm}p{3.5cm}@{}}
  \toprule
  \textbf{Stage} & \textbf{Input Contract} & \textbf{Output Contract} & \textbf{Quality Invariant} \\
  \midrule
  \textbf{Problem Def}  & Business requirements & Measurable objectives; deployment paradigm & All criteria quantifiable \\
  \textbf{Data}          & Objectives; target    & Versioned dataset; validation rules        & Distribution matches production \\
  \textbf{Model Dev}     & Dataset; constraints  & Trained weights; experiment logs            & Meets accuracy within budget \\
  \textbf{Evaluation}    & Model; test data      & Metrics across subgroups                   & No subgroup below threshold \\
  \textbf{Deployment}    & Validated model; SLA  & Serving endpoint; monitoring               & Latency and throughput met \\
  \textbf{Monitoring}    & Live system; baselines & Drift alerts; retraining triggers          & Degradation caught early \\
  \bottomrule
\end{tabular}

\vspace{0.15cm}
\small
\mlsysalert{Violation:}{A missing deployment paradigm at Stage 1 costs $2^{5-1} = 16\times$ to fix at Stage 5.}

\end{frame}

% =============================================================================
\section{Eval \& Deploy}
% =============================================================================

\begin{frame}{Stage 4: Evaluation and Validation}
\note{[3 min] The DR model gets AUC 0.99 on curated data, then sensitivity
drops to 78\% at a rural clinic with older equipment. Lab success does NOT
guarantee production value. Walk through the gap between offline and online
evaluation. Ask: ``Why does a model that works in the lab fail in the field?''}

\small
\begin{columns}[T]
  \begin{column}{0.52\textwidth}
    \textbf{The Lab-to-Field Gap:}
    \begin{itemize}\setlength\itemsep{0pt}
      \item AUC 0.99 on curated research data
      \item Sensitivity \textbf{drops to 78\%} at rural clinic
      \item Older camera, untrained technician
      \item Blurry, poorly-lit, inconsistent framing
    \end{itemize}

    \vspace{0.1cm}
    \textbf{Progressive validation stages:}\\
    {\footnotesize
    1. Offline evaluation $\to$ algorithmic issues\\
    2. Shadow mode $\to$ integration issues\\
    3. Canary (1--5\% traffic) $\to$ scaling issues\\
    4. A/B testing $\to$ user-facing issues}
  \end{column}
  \begin{column}{0.44\textwidth}
    \begin{mlsyscard}{errorstroke}
    {\footnotesize \textbf{Lab $\neq$ Production}\\[0.1cm]
    Medical AI systems experience 10--30\% performance drops in real-world settings. The FDA now requires real-world performance studies.}
    \end{mlsyscard}

    \vspace{0.15cm}
    \mlsysconcept{Key:}{Evaluation measures algorithm quality. Validation confirms production readiness.}
  \end{column}
\end{columns}

\end{frame}

\begin{frame}{Stage 5: Deployment Economics}
\note{[3 min] Quantitative comparison: cloud vs edge for 500 clinics.
Cloud: \$136K/year OpEx. Edge: \$250K CapEx + \$25K/year. Edge pays back
in ~1 year and provides better reliability. But: requires tighter model
optimization. The deployment paradigm selected at Stage 1 determines viability.}

\footnotesize
\textbf{Cloud vs.\ Edge: 500 Clinics, 50 patients/day}

\vspace{0.05cm}
\scriptsize
\renewcommand{\arraystretch}{1.05}
\begin{tabular}{@{}lll@{}}
  \toprule
  & \textbf{Cloud (AWS)} & \textbf{Edge (Jetson)} \\
  \midrule
  Inference cost  & \$0.01/image       & \$0.001/image (electricity) \\
  Annual OpEx     & \textbf{\$136,000} & \$25,000 maintenance \\
  CapEx           & ---                & \$250,000 (one-time) \\
  Latency         & 200+ ms            & \textbf{<50 ms} \\
  Offline?        & No                 & \textbf{Yes} \\
  \midrule
  \textbf{Payback} & --- & \textbf{$\sim$2 years} \\
  \bottomrule
\end{tabular}

\vspace{0.1cm}
\footnotesize
\mlsysconcept{Trade-off:}{Edge: better reliability/latency but tighter model optimization and complex updates.}

\end{frame}

% --- ACTIVE LEARNING 2: Discussion ---
\begin{frame}{Discussion: Which Stage Fails First?}
\note{[3 min] Turn-and-talk. Students discuss in pairs for 90 seconds.
Cold-call 2--3 pairs. The point: all stages are interdependent. A startup
that skips problem definition discovers constraints at deployment.
No single right answer --- the point is systems thinking.}

\centering
\vspace{0.5cm}
{\large\bfseries Turn and Talk \textcolor{midgray}{(90 seconds)}}

\vspace{0.4cm}
{\large A startup deploys a DR screening model.\\
Within 6 months, accuracy drops 15\%.\\[0.2cm]
\alert{Which lifecycle stage failed?}}

\vspace{0.4cm}
\begin{columns}[c]
  \begin{column}{0.16\textwidth}\centering\small Problem\\ Def.\end{column}
  \begin{column}{0.16\textwidth}\centering\small Data\\ Collection\end{column}
  \begin{column}{0.16\textwidth}\centering\small Model\\ Dev.\end{column}
  \begin{column}{0.16\textwidth}\centering\small Evaluation\end{column}
  \begin{column}{0.16\textwidth}\centering\small Deployment\end{column}
  \begin{column}{0.16\textwidth}\centering\small \textcolor{errorstroke}{Monitoring}\end{column}
\end{columns}

\end{frame}

% =============================================================================
\section{Monitoring}
% =============================================================================

\begin{frame}{Stage 6: Silent Degradation}
\note{[3 min] The DR system launches. Six months later, a clinic upgrades
its cameras. Sensitivity drops 8\%. No code changed. No one made an error.
The data drifted. This is the structural difference from traditional
software. Ask: ``How would you detect this without ground truth labels?''}

\footnotesize
\begin{columns}[T]
  \begin{column}{0.52\textwidth}
    \textbf{Six months after launch:}
    \begin{itemize}\setlength\itemsep{0pt}
      \item Clinic upgrades fundus cameras
      \item Sharper images, different color profiles
      \item Sensitivity drops \textbf{8\%} at that site
      \item No code changed. No errors.
      \item \textcolor{errorstroke}{\textbf{Data drifted.}}
    \end{itemize}

    \vspace{0.05cm}
    \textbf{Monitoring hierarchy:}\\
    {\scriptsize
    Operational (seconds) $\to$ Proxy (hours) $\to$ Performance (weeks) $\to$ Drift (months)}
  \end{column}
  \begin{column}{0.44\textwidth}
    \begin{mlsyscard}{errorstroke}
    {\scriptsize \textbf{Production thresholds:}\\
    Sensitivity: >90\% (alert at 88\%)\\
    Specificity: >80\% (alert at 78\%)\\
    Latency P95: <50 ms (alert at 100 ms)\\
    PSI drift: alert if >0.2}
    \end{mlsyscard}

    \vspace{0.05cm}
    \mlsysconcept{Rule:}{ML systems degrade \emph{silently} through data drift.}
  \end{column}
\end{columns}

\end{frame}

\begin{frame}{Multi-Scale Feedback Loops}
\note{[2 min] Five timescales of feedback in the DR system. Each catches
different failures. Fast loops catch operational issues (minutes). Slow
loops catch strategic issues (quarters). The temporal structure prevents
both over-reaction and sluggish adaptation.}

\scriptsize
\renewcommand{\arraystretch}{1.1}
\begin{tabular}{@{}llp{5.5cm}@{}}
  \toprule
  \textbf{Timescale} & \textbf{Loop} & \textbf{What It Catches} \\
  \midrule
  \textcolor{errorstroke}{\textbf{Minutes}}  & Real-time inference  & Misconfigured camera before a day of bad images \\
  \textcolor{routingstroke}{\textbf{Days}}    & Batch monitoring     & Accuracy drift at a specific clinic site \\
  \textcolor{computestroke}{\textbf{Weeks}}   & Drift detection      & Equipment upgrade changed pixel distributions \\
  \textcolor{datastroke}{\textbf{Months}}     & Trend analysis       & Demographic shift requires expanded training data \\
  \textbf{Quarters}   & Strategic review    & Architecture no longer meets clinical needs \\
  \bottomrule
\end{tabular}

\vspace{0.1cm}
\small
\mlsysinsight{Design principle:}{Fast iteration is a \emph{systems feature}, not a convenience.}

\end{frame}

% =============================================================================
\section{Systems Thinking}
% =============================================================================

\begin{frame}{Three Patterns of ML Systems Thinking}
\note{[3 min] Overview of the three structural patterns that recur throughout
the DR case study. These patterns transform reactive debugging into proactive
design. Each pattern appeared at every lifecycle stage.}

% --- Layout: FULL-WIDTH diagram ---
\centering
\safeimg[width=0.88\textwidth,height=4.2cm,keepaspectratio]{ch03-systems-thinking-triad.pdf}

\vspace{0.1cm}
\small
\mlsysconcept{Transformation:}{Reactive debugging $\to$ Proactive design.}

\end{frame}

\begin{frame}{Constraint Propagation Principle}
\note{[3 min] Formalize the Day-150 disaster. Cost grows as 2\^{}(N-1).
A Stage-5 discovery costs 16x. A Stage-6 discovery costs 32x.
The deployment paradigm is the Day-1 constraint, not the last step.}

% --- Layout: FULL-WIDTH chart ---
\centering
\safeimg[width=0.88\textwidth,height=4.0cm,keepaspectratio]{ch03-constraint-propagation.pdf}

\vspace{0.1cm}
\small
\begin{mlsyscard}{crimson}
$\text{Correction Cost} \approx 2^{(N-1)} \times \text{Base Effort}$\\[0.1cm]
{\footnotesize Deployment constraints must flow \textbf{backward} to Day 1. The deployment environment is the \textbf{Day-1 constraint}, not the last step.}
\end{mlsyscard}

\end{frame}

% --- ACTIVE LEARNING 3: Exercise ---
\begin{frame}{Your Turn: Calculate the Cost}
\note{[3 min] Give students 90 seconds. Answer: 2\^{}(6-1) = 32x multiplier.
Must revisit: Data Collection (expand demographics), Model Development
(retrain), Evaluation (re-validate), Deployment (update). That is 4 stages
of rework at exponential cost.}

\small
\begin{columns}[T]
  \begin{column}{0.55\textwidth}
    {\normalsize\bfseries Calculate the Cost}

    \vspace{0.15cm}
    A team discovers during \textbf{Monitoring (Stage 6)} that their DR model fails for patients over 70.

    \vspace{0.1cm}
    This demographic requirement should have been specified at \textbf{Problem Definition (Stage 1)}.

    \vspace{0.15cm}
    \textbf{1.} Calculate the cost multiplier using:\\
    \hspace{1em}$\text{Cost} = 2^{(N-1)}$

    \vspace{0.1cm}
    \textbf{2.} Which stages must be revisited?

    {\footnotesize\textcolor{midgray}{(90 seconds --- then compare with a neighbor)}}
  \end{column}
  \begin{column}{0.42\textwidth}
    \pause
    \begin{mlsyscard}{errorstroke}
    \textbf{Solution:}\\[0.1cm]
    {\footnotesize
    $2^{(6-1)} = 2^5 = $ \textbf{32$\times$} base effort\\[0.1cm]
    \textbf{Stages to revisit:}\\
    2. Data Collection (expand demographics)\\
    3. Model Development (retrain)\\
    4. Evaluation (re-validate)\\
    5. Deployment (update models)\\[0.1cm]
    \textcolor{errorstroke}{\textbf{4 stages}} of rework at exponential cost.}
    \end{mlsyscard}
  \end{column}
\end{columns}

\end{frame}

\begin{frame}{Emergent Complexity}
\note{[2 min] System-level behavior diverges from component-level behavior.
Individual clinics show stable performance, but system-wide analysis detects
subtle degradation across demographic groups. A 2\% accuracy improvement
requires doubling model size, which at 500 clinics means massive CapEx.}

\small
\begin{columns}[T]
  \begin{column}{0.52\textwidth}
    \textbf{Component vs.\ System behavior:}
    \begin{itemize}\setlength\itemsep{0pt}
      \item Individual clinics: stable performance
      \item System-wide: subtle demographic degradation
      \item Invisible in single-site monitoring
    \end{itemize}

    \vspace{0.15cm}
    \textbf{Non-linear trade-offs:}\\
    {\footnotesize
    2\% accuracy gain $\to$ 2$\times$ model size\\
    $\to$ more powerful edge hardware\\
    $\to$ at 500 clinics: massive CapEx}
  \end{column}
  \begin{column}{0.44\textwidth}
    \begin{mlsyscard}{routingstroke}
    {\footnotesize \textbf{ML vs.\ Traditional:}\\[0.1cm]
    Traditional systems fail via \emph{deterministic} cascades (crashes).\\[0.1cm]
    ML systems fail via \emph{probabilistic} degradation (drift, bias amplification).\\[0.1cm]
    Probabilistic degradation lacks obvious error signals.}
    \end{mlsyscard}
  \end{column}
\end{columns}

\end{frame}

% =============================================================================
\section{Wrap-Up}
% =============================================================================

\begin{frame}{Fallacies}
\note{[2 min] Four common misconceptions with quantitative evidence.
Focus on the first two if short on time.}

\footnotesize
\textbf{Fallacy:} \textit{ML development can follow traditional software workflows.}\\
{\scriptsize 60--80\% of ML projects never reach production. Projects forced into waterfall miss the 4--8 iteration cycles production requires.}

\vspace{0.05cm}
\textbf{Fallacy:} \textit{Passing model evaluation means ready for deployment.}\\
{\scriptsize DR model: AUC 0.99 in lab, sensitivity 78\% in rural clinic. Stage-5 discovery costs $16\times$.}

\vspace{0.05cm}
\textbf{Fallacy:} \textit{More data always improves performance.}\\
{\scriptsize Doubling from 500K to 1M often improves accuracy <1\% while doubling costs. 100K clean beats 1M noisy.}

\vspace{0.05cm}
\textbf{Fallacy:} \textit{Data preparation is a one-time step.}\\
{\scriptsize Models degrade 5--10\% within months. Emergency retraining costs 3--5$\times$ proactive monitoring.}

\end{frame}

\begin{frame}{Pitfalls}
\note{[2 min] Three operational pitfalls. If short: cover just the first two.}

\footnotesize
\textbf{Pitfall:} \textit{Skipping validation to accelerate timelines.}\\
{\scriptsize Skipping shadow mode causes 10--50$\times$ latency spikes. ``Saving'' 2 weeks costs 6--8 weeks of emergency remediation.}

\vspace{0.08cm}
\textbf{Pitfall:} \textit{Deferring deployment paradigm selection until after model development.}\\
{\scriptsize A team develops a 2 GB ensemble, discovers TinyML target with 256 KB. Constraint Propagation: $2^4 = 16\times$ rework.}

\vspace{0.08cm}
\textbf{Pitfall:} \textit{Treating the model as the sole deliverable.}\\
{\scriptsize Reproducible artifact = weights + code + environment + config. Without bundling, ``works on my machine'' = broken.}

\end{frame}

% --- RETRIEVAL PRACTICE ---
\begin{frame}{What Were the Key Ideas?}
\note{[2 min] Retrieval practice. Students write for 90 seconds, no notes.
Do NOT show next slide yet. Walk around the room.}

\centering
\vspace{1.5cm}
{\Large\bfseries Close your notes.}

\vspace{0.8cm}
{\large Write down the \textbf{4 most important concepts} from today.}

\vspace{0.8cm}
{\normalsize\textcolor{midgray}{90 seconds --- no peeking.}}

\end{frame}

\begin{frame}{Key Takeaways}
\note{[2 min] Reveal. Walk through each bullet. Emphasize the three
quantitative anchors: 60--80\% data time, 4--8 iteration cycles, 2\^{}(N-1)
cost escalation.}

\footnotesize
\begin{itemize}\setlength\itemsep{1pt}
  \item \textbf{Two pipelines, one system}: Data and model development run in parallel, unified by continuous feedback loops.
  \item \textbf{60--80\% of effort is data-related}: Model development is only 10--20\%. Plan accordingly.
  \item \textbf{Iteration velocity compounds}: 10 experiments/day beats 1/week, even starting 5\% behind (Iteration Tax).
  \item \textbf{Constraint propagation is exponential}: $\text{Cost} = 2^{(N-1)}$. A Stage-5 discovery costs $16\times$ base effort.
  \item \textbf{Each stage maps to an Iron Law term}: Data $\to D_{\text{vol}}$, Model $\to O$, Deploy $\to L_{\text{lat}}$.
  \item \textbf{ML systems degrade silently}: Data drift, not code bugs. Continuous monitoring from day one.
  \item \textbf{Systems thinking}: Constraint propagation + multi-scale feedback + emergent complexity.
\end{itemize}

\end{frame}

\begin{frame}{References}
\note{[1 min] Point students to canonical papers.}

\small
\mlsysref{Sculley+15}{Sculley et al. ``Hidden Technical Debt in ML Systems.'' NeurIPS 2015.}
\mlsysref{CrowdFlower16}{CrowdFlower. ``Data Science Report.'' 2016.}
\mlsysref{Amershi+19}{Amershi et al. ``Software Engineering for ML.'' ICSE 2019.}
\mlsysref{Gulshan+16}{Gulshan et al. ``Development of a Deep Learning Algorithm for Detection of Diabetic Retinopathy.'' JAMA 2016.}
\mlsysref{Paleyes+22}{Paleyes et al. ``Challenges in Deploying ML.'' ACM Comp.\ Surveys 2022.}
\mlsysref{Boehm81}{B. Boehm. \textit{Software Engineering Economics}. 1981.}

\end{frame}

\begin{frame}{Next Lecture: Data Engineering}
\note{[1 min] Forward hook. The workflow framework provides the map. Now we
need fuel: data. Data engineering is where 60--80\% of the effort goes, and
where most failures begin. The question: how do you build pipelines that
feed the right data to the right model at the right time?}

\footnotesize
\begin{columns}[c]
  \begin{column}{0.30\textwidth}
    \centering
    \begin{mlsyscard}{datastroke}
    {\large\bfseries Data Pipelines}\\[0.1cm]
    {\footnotesize Collection, cleaning, versioning}
    \end{mlsyscard}
  \end{column}
  \begin{column}{0.30\textwidth}
    \centering
    \begin{mlsyscard}{computestroke}
    {\large\bfseries Feature Eng.}\\[0.1cm]
    {\footnotesize Transform raw data into ML-ready features}
    \end{mlsyscard}
  \end{column}
  \begin{column}{0.30\textwidth}
    \centering
    \begin{mlsyscard}{errorstroke}
    {\large\bfseries Data Quality}\\[0.1cm]
    {\footnotesize Validation, drift detection, monitoring}
    \end{mlsyscard}
  \end{column}
\end{columns}

\vspace{0.3cm}
\centering
{\normalsize The engine needs fuel.}\\[0.1cm]
{\small\textbf{60--80\% of ML effort is data engineering. That is where we go next.}}

\end{frame}

\end{document}
